% FILE: 07_open_questions_and_next_work.tex
% Project: standards
% Thesis Role: public-standards-compliance-and-gravity-field-mapping

\section{Open Questions and Next Work}
\label{sec:standards-open}

\subsection{Known Gaps}
\label{sec:standards-open-gaps}

Four structural gaps are documented as of the ALIVE checkpoint. These are not gaps in the
completed crosswalk---all 16 verified standards remain ALIVE---but gaps in the boundary
between the crosswalk and the downstream execution surface.

\paragraph{OCPQ query evaluation engine (AUTHOR THESIS).}
The OCPQ placement file is VERIFIED in the audit matrix. Structural shapes for object-centric
process queries exist in \texttt{src/ocpq.rs} in the wasm4pm codebase. However, the query
evaluation engine itself has not graduated to wasm4pm: the shapes define the query language
surface but do not yet execute queries against OCEL 2.0 logs. This is documented as a
graduation boundary item. Until the evaluation engine graduates, OCPQ claims about query
answering are PARTIAL at the execution level, though ALIVE at the type-law level.

\paragraph{SHACL shapes for OCEL 2.0 RDF representation.}
The SHACL placement file describes the Shapes Constraint Language validation surface for
OCEL 2.0 RDF representations. No \texttt{.shacl} or \texttt{.ttl} files exist in the
\texttt{standards/} directory. The surface is documented but not materialized as a
machine-executable artifact. This means SHACL validation of OCEL 2.0 RDF cannot currently
be run without authoring the shape files.

\paragraph{OTel Weaver process-specific semantic conventions.}
The OTel Weaver placement maps the Weaver framework to the process intelligence context.
However, the actual \texttt{.tera} template files and \texttt{.yaml} convention files that
would define process-specific semantic conventions (e.g., activity name attribution, lifecycle
state encoding, object type labeling) are not present in the \texttt{standards/} directory.
The convention framework is mapped; the conventions themselves are deferred.

\paragraph{Log Skeleton (Verbeek 2018) placement.}
Log Skeleton is noted in the project materials as having structural shape coverage via
\texttt{DeclareConstraint}, but it does not appear among the 16 verified standards and has
no dedicated placement file in the audit matrix. Its conformance checking capability has
not graduated.

\subsection{Deferred Gates}
\label{sec:standards-open-deferred}

The following proof gates are deferred to future work phases:

\begin{enumerate}
  \item \textbf{OCPQ execution graduation}: A gate requiring that the OCPQ evaluation
        engine execute at least one multi-object query against a production-grade OCEL 2.0
        log and emit a receipt recording the query AST, the result set hash, and the
        non-forgeability witness.
  \item \textbf{SHACL/OCEL RDF validation}: A gate requiring that SHACL shape files for
        OCEL 2.0 RDF representation be authored, published, and validated against at least
        one conformant and one non-conformant OCEL 2.0 RDF graph, with deterministic
        validation receipts.
  \item \textbf{OTel Weaver convention authoring}: A gate requiring that at least three
        process-specific semantic conventions be authored in Weaver \texttt{.yaml} format
        and validated against the OTel trace integration schema.
  \item \textbf{Empirical competitor conformance measurement}: The Nash equilibrium cage
        analysis in \texttt{reverse-lock-in.md} describes expected competitor behavior but
        records no executed experiment with a measured competitor conformance adoption rate.
        A future gate requires at least one documented case of a competing system's traces
        being evaluated against $W_{\text{std}}$ with a recorded fitness score.
\end{enumerate}

\subsection{Research Risks}
\label{sec:standards-open-risks}

\paragraph{Nash equilibrium empirical gap.}
The reverse lock-in mathematical formalism is a strategic doctrine supported by formal
game theory and competitor payoff analysis. However, no executed experiment or measured
competitor conformance adoption rate is recorded. The Game-Theoretic Equilibrium section
of \texttt{reverse-lock-in.md} describes expected behavior. This is documented as a known
gap; until empirical data exists, the Nash equilibrium claim is INTERPRETATION / FUTURE
WORK rather than a finding supported by evidence.

\paragraph{OCPQ expressiveness boundary.}
The OCPQ placement documents query AST caching and result-set non-forgeability, but the
boundary of the OCPQ query language with respect to OCEL 2.0 expressiveness has not been
formally characterized. It is possible that some multi-object queries required by the
M\&A evidence path fall outside the current OCPQ language definition.

\paragraph{SHACL/OCEL 2.0 shape completeness.}
The OCEL 2.0 bipartite graph invariants are formally defined in the placement file, but
their translation into SHACL shapes may encounter RDF serialization issues. The global
Python/RDF rules note that the \texttt{format:turtle} serializer has known issues; OCEL
2.0 RDF shapes will need to be authored in N-Triples format and validated accordingly.

\subsection{Next Receipted Work}
\label{sec:standards-open-next}

The next receipted work items, each of which will produce a named checkpoint on the ledger,
are:

\begin{enumerate}
  \item Author SHACL shape files for OCEL 2.0 RDF and issue a validation receipt recording
        at least one conformant and one non-conformant result.
  \item Graduate the OCPQ query evaluation engine to wasm4pm and issue a query execution
        receipt recording the query AST hash, result set size, and non-forgeability witness.
  \item Author OTel Weaver \texttt{.yaml} semantic conventions for process intelligence
        (minimum: activity name, lifecycle state, object type) and validate against the
        existing OTel trace integration schema.
  \item Design and execute an empirical measurement of the reverse lock-in doctrine: apply
        $W_{\text{std}}$ conformance checking to at least one external process model and
        record the fitness score, precision, and any Blue River Dam rejection events.
  \item Author a Log Skeleton (Verbeek 2018) placement file and graduate it as standard 17
        in the audit matrix, with a signed checkpoint update.
\end{enumerate}

Each of these items will close a documented gap and extend the ALIVE scope of the standards
project. None may be claimed ALIVE until the receipt, tests, and checkpoint all exist per
the foundry's immutability doctrine.
